\chapter*{Conclusiones y recomendaciones}
\addcontentsline{toc}{chapter}{Conclusiones y recomendaciones}
\markboth{Conclusiones y recomendaciones}{Conclusiones y recomendaciones}

\section*{Conclusiones}
\phantomsection\addcontentsline{toc}{section}{Conclusiones}

El trabajo cumplió el objetivo general: se desarrolló EduFEM, un software educativo de elementos finitos que expone el procedimiento de cálculo de problemas bidimensionales de elasticidad lineal, en tensión plana y en deformación plana, con elementos cuadriláteros isoparamétricos Q4 y Q9. Se diseñó e implementó en el lenguaje de programación Python, y se verificó y validó frente a soluciones analíticas, casos clásicos de la literatura y un software comercial. La finalidad que el objetivo enuncia, que el análisis por el MEF resulte trazable y verificable paso a paso, queda establecida sobre el software y sus resultados con los criterios de la \autoref{sec:validacion-resultados}; su efecto en la enseñanza de la Carrera no forma parte de lo medido, según se precisa más adelante.

El problema científico preguntaba de qué manera el uso del software de análisis estructural como caja negra podrá incidir en la baja trazabilidad y verificabilidad del análisis por el Método de los Elementos Finitos en elasticidad plana, en la Carrera de Ingeniería Civil de la Universidad Autónoma Tomás Frías, en la gestión 2026. La respuesta que este trabajo sostiene es de orden técnico: la incidencia está en el medio de cálculo, que es el campo de acción declarado en la Introducción. Un software que solo muestra datos y resultados deja un análisis que no puede seguirse hasta sus datos ni comprobarse paso a paso; uno que expone el procedimiento permite ambas cosas. El trabajo lo establece por la vía que le corresponde a una investigación tecnológica: construyendo ese software y midiendo sobre él las dos propiedades, con criterios fijados de antemano.

El alcance de esa respuesta determina cómo debe leerse todo lo que sigue. La evidencia procede de evaluar el software y sus resultados en escenarios controlados. No procede de observar su uso en la Carrera ni de medir efecto alguno sobre quienes lo emplean: la Carrera es el contexto al que el trabajo se destina y en el que delimita su problema, no un lugar del que se hayan tomado observaciones. Queda establecido, con los criterios fijados de antemano, que un medio de cálculo puede tener las dos propiedades y que EduFEM las tiene: la verificabilidad, con evidencia numérica contrastada contra patrones independientes; la trazabilidad, con el inventario y el cotejo documental que describe la \autoref{sec:validacion-datos}. Queda fuera de lo medido qué ocurre cuando ese medio se pone en manos de un curso. Las conclusiones específicas siguen el orden de los objetivos.

\textbf{Objetivo específico 1 —diagnosticar—} (\autoref{sec:antecedentes}). La revisión de los antecedentes y el análisis comparativo de la \autoref{tab:comparativa} establecieron que ninguno de los recursos revisados documenta los atributos que hacen trazable y verificable un análisis por el MEF en elasticidad plana. El software comercial oculta el procedimiento por diseño. El antecedente educativo más directo, ED-Elas2D, lo expone en parte, pero sus propios autores describen su post-proceso como una herramienta meramente descriptiva \autocite[p.~253]{suarez1998edelas2d}. De ese diagnóstico salieron los cinco requisitos de la herramienta, enumerados al cierre de la \autoref{sec:antecedentes} y desarrollados en la \autoref{sec:requisitos}:
\begin{itemize}
    \item el procedimiento de cálculo completo a la vista sobre el modelo del usuario;
    \item una memoria de cálculo reproducible a mano;
    \item una exactitud verificada y validada, publicada junto con la herramienta;
    \item la cobertura del contenido que exige un consenso publicado de expertos;
    \item un software libre y en español.
\end{itemize}

\textbf{Objetivo específico 2 —fundamentar—} (\autoref{cap:marco_teorico}). Se estableció la base que comparten el motor, los módulos educativos y la memoria de cálculo: la formulación del MEF para la elasticidad lineal plana con elementos Q4 y Q9, de la formulación variacional a la recuperación de tensiones; la calidad de malla; los fenómenos numéricos; y los criterios de verificación y validación. Se fijaron también los cuatro principios que orientan la exposición del procedimiento: su visibilidad, la interactividad con vínculo al modelo, la construcción paso a paso y la conexión entre el fundamento y la operación (\autoref{sec:fundamentos-pedagogicos}). Esa base no es un resultado medido, sino el fundamento del trabajo. Lo que los resultados comprueban es que el motor construido sobre ella se comporta como la teoría predice, y que las expresiones del capítulo son las que el código implementa, los módulos exhiben y la memoria sustituye numéricamente.

\textbf{Objetivo específico 3 —desarrollar—} (\autoref{sec:diseno-edufem}). Se construyó un motor de cálculo que cubre las nueve etapas del procedimiento para elementos isoparamétricos: del mapeo isoparamétrico y el Jacobiano a la imposición de restricciones —incluidos los desplazamientos prescritos no nulos— y la recuperación de tensiones por extrapolación desde los puntos de Gauss. Sobre él se construyó una interfaz organizada en pre-proceso, proceso y post-proceso, que comparten un único lienzo, con selección en ambos sentidos entre las tablas y el lienzo, deshacer y rehacer sobre el modelo completo, y comprobación de la salud del modelo antes de resolver. Ocho módulos educativos exponen las etapas sobre el modelo real del usuario: del mapeo isoparamétrico al ensamblaje (M1 a M7, \autoref{tab:modulos}), más la calidad de malla en el pre-proceso (M0). El post-proceso expone la solución y la recuperación de tensiones en modo crudo y suavizado, con los medios de consulta de la \autoref{sec:post-proceso}, entre ellos la sonda con su círculo de Mohr y la vista tridimensional. La memoria de cálculo en PDF reconstruye el procedimiento paso a paso con las mismas funciones del motor. Los formatos de intercambio completan el recorrido fuera de la herramienta: el modelo exportado al paquete CSV/ZIP vuelve a ella sin pérdida de entidades, y reimportar un DXF no crea entidades nuevas (\autoref{sec:consistencia-interna}).

\textbf{Objetivo específico 4 —verificar y validar—} (\autoref{sec:mms} a \autoref{sec:cook}). Se aplicó una batería de tres casos de estudio, contra los criterios de aceptación fijados de antemano en la \autoref{sec:validacion-resultados} y codificados en los guiones de verificación y validación. La verificación por el método de soluciones manufacturadas, en los dos estados planos y sobre mallas regulares y distorsionadas, arrojó las tasas de convergencia que la teoría predice para el desplazamiento, es decir, el ritmo al que el error disminuye al refinar la malla. Esas tasas son $2{,}00$ en la norma $L^2$ del desplazamiento y $1{,}00$ en la seminorma $H^1$ de su gradiente para el Q4, y $3{,}00$ y $2{,}00$ para el Q9 (\autoref{tab:mms}). Esas tasas comprueban que, en las cuatro configuraciones ensayadas, el ensamblaje, la cuadratura, el mapeo isoparamétrico, la matriz constitutiva y la imposición de restricciones no introducen errores que degraden el orden de convergencia. La corrección de la extrapolación de tensiones se comprueba aparte, con la reproducción exacta de campos polinómicos. En la viga de Timoshenko, los errores frente a la solución analítica fueron del $0{,}04\,\%$ en la tensión normal $\sigma_x$ y del $0{,}26\,\%$ en la flecha; la diferencia frente al modelo de SAP2000 fue de hasta el $0{,}21\,\%$ en $\sigma_x$. Esos máximos se toman sobre los tres puntos de control; en las componentes secundarias, de magnitud entre seis y treinta y cinco veces menor, el error llega al $2{,}89\,\%$ frente a la solución analítica (\autoref{tab:timoshenko-stress}). La suma de las reacciones verticales equilibra la carga total con un residuo relativo de $1{,}7\times10^{-13}$. La membrana de Cook permitió ilustrar numéricamente el bloqueo por cortante (\emph{shear locking}) del elemento Q4 frente a la convergencia rápida del Q9: con 578 grados de libertad, el error del Q9 es más de diez veces menor que el del Q4. Esta validación se apoya en soluciones analíticas y en un modelo equivalente en un software comercial, según la acepción declarada en la \autoref{sec:verificacion-validacion}; no se contrastó contra ensayos físicos.

\textbf{Objetivo específico 5 —validar la propuesta—} (\autoref{sec:cobertura} a \autoref{sec:validacion-hipotesis}). Los criterios de trazabilidad se cumplen. Siete de las siete primeras etapas tienen módulo educativo. Las nueve etapas tienen desarrollo con sustitución numérica en la memoria de cálculo. Las tres fases operan sobre un mismo lienzo. Y el autor cotejó, ítem por ítem y con la regla de marcado de la \autoref{sec:validacion-datos}, que los dieciocho ítems del universo de contenido tienen instrumento en EduFEM: las quince destrezas y los tres conceptos que un consenso publicado de 67 expertos exige dentro del alcance del trabajo \autocite[pp.~1162-1163]{perezsantiago2023fem}. La memoria del ejemplo canónico se desarrolla paso a paso en el \autoref{anexo:memoria}, con cada valor remitido a la ecuación del \autoref{cap:marco_teorico} de la que procede, para que pueda rehacerse a mano. La \autoref{tab:veredicto} reúne, criterio por criterio, el umbral, la cifra observada y el veredicto.

En conjunto, el trabajo cumple los cinco objetivos específicos y, con ellos, el objetivo general. Los dos primeros se cumplen como diagnóstico y como fundamento, no como resultados medidos. El tercero se cumple con módulos educativos hasta el ensamblaje, y con las dos últimas etapas expuestas en el post-proceso y en la memoria.

\textbf{La hipótesis queda comprobada en sus dos cláusulas, sobre el software y los casos de estudio.} En la cláusula (a), trazabilidad, todos los criterios de cobertura se cumplen. En la cláusula (b), verificabilidad, todos los criterios numéricos se cumplen dentro de las tolerancias fijadas de antemano. El desarrollo de EduFEM mejoró la forma en que el software expone el procedimiento de cálculo —de la caja negra que caracteriza la \autoref{tab:comparativa} al procedimiento a la vista— y, con ello, el análisis por el MEF en elasticidad plana resulta trazable y verificable, dentro del alcance declarado. El propio desarrollo aportó una comprobación no prevista de esa relación. Un defecto en la extrapolación de tensiones del elemento Q4 no lo delataron los estudios de convergencia del desplazamiento, porque esa matriz no interviene en el desplazamiento; lo delató la lectura de la memoria de cálculo, donde las tensiones nodales no cerraban con las de los puntos de Gauss (\autoref{sec:leccion-extrapolacion}). En un software de caja negra, ese defecto habría sido invisible para quien lo usa.

\section*{Aportes del trabajo}
\phantomsection\addcontentsline{toc}{section}{Aportes del trabajo}

Una investigación en ciencia del diseño se distingue de un diseño rutinario por lo que deja en la base de conocimiento \autocite[p.~81]{hevner2004design}. Este trabajo deja tres aportes.

El primero es el software. EduFEM es una herramienta libre, en español y verificada, que reúne tres capas que los recursos revisados como antecedentes presentan por separado: un motor de cálculo, una interfaz de modelado interactiva y un conjunto de módulos educativos. El motor implementa la formulación isoparamétrica Q4 y Q9 para los dos estados planos de la elasticidad, verificada por soluciones manufacturadas en ambos estados y contrastada con dos referencias externas en tensión plana. El contraste con ambas referencias sobre la viga de Timoshenko deja el error de la tensión normal $\sigma_x$ y de la flecha por debajo del $0{,}3\,\%$ frente a la solución analítica, y la diferencia en $\sigma_x$ frente a SAP2000, por debajo del $0{,}3\,\%$.

El segundo es conocimiento de diseño: tres decisiones que hacen trazable un software de elementos finitos y que pueden reutilizarse. La primera es que cada módulo educativo opere sobre el elemento que el usuario selecciona en el lienzo compartido, y no sobre un ejemplo prefabricado; así, fenómenos como el bloqueo por cortante o los modos espurios (\emph{hourglass}) pueden observarse en el modelo propio. La segunda es la memoria de cálculo generada con las mismas funciones del motor: un documento que reconstruye el procedimiento con sustitución numérica, de las funciones de forma a las tensiones nodales, y lo cierra con un diagnóstico de equilibrio y de condicionamiento. El estudiante puede seguirlo, rehacerlo por sus propios medios y usarlo para comprobar los resultados. La tercera es conservar, junto al motor rápido, una versión legible del cálculo elemento por elemento, que sirve a la vez de referencia para los módulos y de patrón de comparación para las pruebas de regresión (\autoref{sec:motor}).

El tercero es conocimiento de evaluación. El trabajo deja una batería de verificación y validación reproducible, con sus guiones y sus datos, y un ejemplo documentado de cómo se verifica un programa de cálculo estructural. Deja también la prueba de reproducción exacta de campos polinómicos, que el límite de la propia batería obligó a incorporar (\autoref{sec:leccion-extrapolacion}). Y deja la tabla de cobertura de contenido, construida sobre los ítems de un consenso publicado de expertos. Se aplicó aquí a un único software y por un único evaluador; su uso con otros recursos queda por ensayar.

\section*{Limitaciones}
\phantomsection\addcontentsline{toc}{section}{Limitaciones}

Las fronteras del trabajo se fijaron de antemano en la Introducción (\emph{Alcance y limitaciones}). A ellas se suman las que el desarrollo y los resultados pusieron de manifiesto, que la \autoref{sec:interpretacion} discute y que aquí se resumen.

\begin{itemize}
    \item La evaluación la hizo el autor. El software fue ejercitado por el autor y por su batería automática de pruebas; no lo operaron terceros, y su facilidad de uso no está establecida. El motor, los guiones de verificación y los umbrales los escribió el autor, sin revisión independiente del código. La contrastación externa se apoya en la solución analítica de Timoshenko y Goodier y en el modelo de SAP2000.
    \item La tabla de cobertura es un cotejo del autor sobre un criterio externo publicado, no el juicio de un panel independiente. Ese criterio procede de expertos mexicanos de ingeniería mecánica, y el propio estudio advierte que sus resultados no son generalizables a otros contextos \autocite[p.~1171]{perezsantiago2023fem}.
    \item La exposición del procedimiento es una posibilidad que el software ofrece, no una obligación que imponga. Un usuario puede construir el modelo, resolverlo y leer los resultados sin abrir ningún módulo ni generar la memoria. La herramienta no fuerza el recorrido, porque el orden de exploración lo fija la enseñanza y no el software.
    \item Los módulos educativos llegan hasta el ensamblaje, por decisión de diseño. La solución del sistema y la recuperación de tensiones se exponen en el post-proceso y en la memoria, y no en un módulo propio. Las destrezas de mallado se ejercitan sobre mallas que el usuario debe construir sin un generador automático.
    \item El elemento Q4 presenta bloqueo por cortante bajo flexión dominante y, en deformación plana, la matriz constitutiva degenera cuando $\nu \to 0{,}5$; el software señala ambos fenómenos, pero no los corrige (\autoref{sec:interpretacion}).
    \item La deformación plana solo se verifica con el método de soluciones manufacturadas: los dos casos de contraste externo son de tensión plana, y la viga de Timoshenko se resolvió únicamente con elementos Q9. La carga superficial linealmente variable no interviene en ningún caso de estudio.
    \item El campo de tensiones recuperado del Q4 converge a $\mathcal{O}(h^{1{,}54})$ y no alcanza $\mathcal{O}(h^{2})$. La explicación que propone el trabajo —la capa de contorno, donde cada nodo recibe el promedio de un solo lado— no se contrastó (\autoref{sec:mms}).
    \item El solucionador es directo (\autoref{sec:motor}). Es adecuado para los tamaños de modelo de un contexto educativo y para las mallas de decenas de miles de grados de libertad que mide la \autoref{tab:tiempos}, en las que el almacenamiento disperso de $\bm{K}$ crece casi en proporción al número de grados de libertad. Su límite está en las mallas muy grandes: el tiempo de solución crece más que linealmente con el número de grados de libertad, porque la factorización se vuelve dominante. La memoria que consume la propia factorización no se midió.
\end{itemize}

\section*{Recomendaciones y trabajo futuro}
\phantomsection\addcontentsline{toc}{section}{Recomendaciones y trabajo futuro}

A partir de las limitaciones anteriores y de las fronteras que la Introducción fijó de antemano se plantean tres líneas de continuación, una por capítulo.

\subsection*{Sobre la formulación (\autoref{cap:marco_teorico})}
\phantomsection\addcontentsline{toc}{subsection}{Sobre la formulación}

\textbf{Tratamiento del bloqueo.} Para abordar el bloqueo (\autoref{sec:locking}) —por cortante en flexión y volumétrico cuando $\nu \to 0{,}5$ en deformación plana— se recomienda incorporar técnicas correctivas, como la integración reducida selectiva o la formulación $\bar{\bm{B}}$ \autocites[pp.~221, 232-234]{hughes2000fem}[pp.~476-477]{bathe2014fem}. El Q9 con integración completa tampoco está libre del bloqueo volumétrico: estas técnicas beneficiarían a ambos elementos. Permitirían, además, un módulo educativo que contraste, sobre un mismo modelo, el comportamiento bloqueado y el corregido: el fenómeno que hoy solo se muestra con la membrana de Cook pasaría a poder manipularse.

\textbf{Más tipos de elemento.} La biblioteca de elementos puede ampliarse con elementos triangulares, lineales y cuadráticos, y con cuadriláteros de ocho nodos y de transición, para enriquecer el estudio comparativo de la convergencia y del comportamiento bajo distorsión de malla \autocites[pp.~167-174, 194]{onate2009structural}{reddy2006introduction}.

\subsection*{Sobre el software (\autoref{cap:diseno})}
\phantomsection\addcontentsline{toc}{subsection}{Sobre el software}

\textbf{Mallado automático.} Un generador de mallas sobre un contorno dibujado en el lienzo, con las métricas de calidad del módulo M0 como control, reduciría el trabajo manual que hoy exige un modelo de más de unas decenas de elementos. Convertiría además las destrezas de mallado de la tabla de cobertura —selección del tamaño, refinamiento en zonas de concentración, prueba de convergencia— en operaciones de la propia herramienta.

\textbf{Solucionador para problemas de gran tamaño.} Las extensiones de mayor impacto sobre el esquema descrito en la \autoref{sec:motor} son dos: una factorización de Cholesky, aplicable porque $\bm{K}_{ff}$ es simétrica y definida positiva (\autoref{sec:resolucion-tensiones}), y un solucionador iterativo precondicionado para mallas muy grandes. Ambas preservarían la versión legible del ensamblaje que usan los módulos educativos.

\textbf{Guía de uso en la cátedra.} El manual del \autoref{anexo:manual} describe la operación del software. Conviene acompañarlo con una guía docente que proponga una secuencia de prácticas. Tres son inmediatas: construir y resolver el ejemplo canónico y rehacer a mano una etapa de su memoria de cálculo; recorrer los módulos M1 a M7 sobre un elemento propio; y convertir un modelo de Q4 a Q9 y refinar la malla sobre la membrana de Cook para observar el bloqueo por cortante.

\textbf{Extensión a tres dimensiones.} A más largo plazo, la extensión del motor a problemas tridimensionales con elementos hexaédricos y tetraédricos sería un paso natural, porque la formulación isoparamétrica y la infraestructura de cuadratura se formulan del mismo modo en tres dimensiones.

\subsection*{Sobre la evidencia (\autoref{cap:resultados})}
\phantomsection\addcontentsline{toc}{subsection}{Sobre la evidencia}

\textbf{Ampliación de la batería de verificación y validación.} Se recomienda incorporar un caso con carga superficial linealmente variable, que ningún caso de estudio ejercita, y al menos un caso civil en deformación plana con solución publicada —una presa de gravedad, un muro de contención o un túnel somero— con peso propio. Así, la validación externa cubriría el estado plano que la práctica civil emplea con más frecuencia. En la verificación queda un experimento pendiente: recomputar la norma $L^2$ del campo de tensiones del Q4 excluyendo la franja de contorno, que contrastaría la explicación de su tasa de 1,54 (\autoref{sec:mms}).

\textbf{Revisión independiente.} Se recomienda someter la tabla de cobertura de contenido al juicio de docentes de la asignatura o de especialistas en el MEF, y el código del motor a una revisión por terceros. Una línea de investigación distinta, de carácter educativo, podrá estudiar el uso de la herramienta en el aula.

De las tres líneas, las dos primeras amplían el software y la tercera amplía la evidencia que lo respalda y deja señalada una línea educativa, ajena al alcance de este trabajo. Las tres conservan lo que define a EduFEM: que cada resultado pueda seguirse hasta sus datos y comprobarse paso a paso.
